\refstepcounter{chapter}
\chapter*{EK A}
Ekler bölümünde verilen şekil ve çizelgeler, bulundukları ekin adı altında numaralandırılabilir. LaTeX varsayılanı bu bölümde A.1, A.2 şeklindedir.

\begin{figure}[htbp]
\centering
\fbox{\rule{0pt}{3cm}\rule{6cm}{0pt}}
\caption{Ekler bölümünde şekil örneği.}
\label{fig:ek-ornek}
\end{figure}

\begin{table}[htbp]
\caption{Ekler bölümünde çizelge örneği.}
\label{tab:ek-ornek}
\centering
\begin{tabular}{c c c c}
\toprule\toprule
Kolon A & Kolon B & Kolon C & Kolon D\\
\midrule
Satır A & Satır A & Satır A & Satır A\\
Satır B & Satır B & Satır B & Satır B\\
Satır C & Satır C & Satır C & Satır C\\
\bottomrule
\end{tabular}
\end{table}
